\section{Discussion}
\label{sec:discussion}

\textbf{Answer to the research question.} With a 0.6B model, no: the agent did not beat ADR in any
scenario, degraded the network in the traffic surge and under heavy load, and its few actions were
almost always the same template rather than a response to the observed event. The event-triggered architecture, however,
is clearly useful: the same triggers, telemetry and leased overrides driven by six hand-written
rules gain 10--15 points in three of the five scenarios. An LLM agent is only worth its cost if it
beats that rule set, for instance on combinations of events that no rule anticipates. Whether a
3B--7B model clears that bar is the next measurement; the scripts run unchanged with Ollama
(\texttt{--backend ollama --model qwen2.5:3b}).

\textbf{Validated is not safe.} Range validation, the read-before-write rule and the tool cap all
held, yet the agent's only recurrent action (every device to SF12) was harmful during a surge, and in
one run its explanation contradicted its action. We tested a predicted-load check (reject an action whose
offered load would exceed the ALOHA optimum $G=0.5$) and found that it would not have blocked this
case, because the surge had not been observed yet when the agent acted. Effective protection needs
to act on consequences: roll an override back when the measured loss rises after it, apply it first
to a subset of devices, or simulate it in a digital twin of the network before applying it --- the
simulator of this work could serve as that twin.

\textbf{Physics first.} The Doppler rate makes the best SF over a LEO pass change within minutes,
much faster than ADR's 20-uplink history, and whether SF12 survives near zenith depends on the
receiver's frequency tracking (Fig.~\ref{fig:e2b}). A controller that knows the ephemeris and the
receiver's Doppler limits can choose the SF from the elevation; ADR cannot.

\textbf{The control plane bounds every controller.} Under class~A, a command waits for the next
delivered uplink. Agent latency (seconds) is negligible against this delay (minutes), and a deep
fade cannot be fixed from the network at all; only device-side logic or class~B/C downlinks could.

\textbf{Oracles and imitation.} A per-packet oracle that minimises energy is not an upper bound under
heavy load, and a classifier that imitates it without its information under-provisions. Labels
should target an expected PER given what the network can observe, and should account for collisions.

\textbf{Limitations.} No Earth rotation; S-band TR~38.811 tables at 868\,MHz; orthogonal SFs; no
background traffic at the satellite; a broadcast satellite downlink is assumed; one pass per
scenario and 3\,h horizons; receivers A and B are models, not the Semtech demodulator; the agent
results come from a 0.6B model run through llama.cpp because the development environment could not
reach Ollama. All assumptions are listed in the repository (\texttt{docs/ASSUMPTIONS.md}).

\section{Conclusion}
We built a reproducible, Doppler-aware simulation of a hybrid terrestrial--LEO LoRa network and
evaluated an event-triggered LLM agent against fixed SFs, ADR, an ML classifier, an oracle and a
rule-based controller on the same triggers. The physical layer reproduces the measured SF12 limit
of LEO links within about 100\,km of altitude. The event-triggered architecture gains about ten
points of delivery ratio during jamming, satellite pass ends and gateway outages. A 0.6B agent did
not exploit it and once caused harm that schema validation could not stop. Next steps: the 3B and
7B models, a consequence-aware safety layer, store-and-forward buffering, and several passes per
day with Earth rotation.
